\chapter{Division rings and real division algebras}
\label{chap:M07.23}
The word ring in this chapter retains associativity and an identity, but
commutativity is not assumed. The underlying division rings are nonzero.
This distinction is essential: the earlier finite-domain argument assumed
commutativity, whereas the finite-division-ring theorem below proves it.

\KEntry{K07.23.centralizers}{Centers, centralizers and dimensions over a division ring}
\begin{definition}
A division ring $D$ is a nonzero associative unital ring in which every
nonzero element has a two-sided multiplicative inverse. Its center is
$Z(D)=\{a:ax=xa\text{ for every }x\in D\}$. For $x\in D$ its
centralizer is $C_D(x)=\{a:ax=xa\}$.
\end{definition}
Both sets are division subrings. Closure under addition and multiplication
follows by distributing the commuting identities. If $ax=xa$ and $a\ne0$,
multiply this identity on both sides by $a^{-1}$ to obtain
$a^{-1}x=xa^{-1}$. The same argument for every $x$ treats the center.
The center is commutative, hence a field. A subring containing $1$ in a
finite division ring is itself a division ring: multiplication by a
nonzero element is injective on the finite subring, hence surjective, and
its inverse therefore lies in that subring.

For later counting, let $E$ be a division subring of a finite division
ring $D$. Regard $D$ as a left $E$-space. Choose a maximal left-linearly
independent list $v_1,\ldots,v_r$, which exists by finiteness. If a vector
$v$ is not in its span, a relation
$a v+\sum a_i v_i=0$ with $a\ne0$ would express $v$ as an old linear
combination after multiplication by $a^{-1}$ on the left. Thus adjoining
$v$ would preserve independence, contradicting maximality. The list
spans and its coordinate expression is unique by independence. Consequently
$|D|=|E|^r$. This argument uses the side of scalar multiplication consistently;
it does not assume that $E$ is commutative.

If $F\subseteq Z(D)$ is a finite field with $q$ elements and $E$ contains
$F$, the same argument over the field $F$ gives $|E|=q^d$ and
$|D|=q^n$, where $n=dr$. In particular $d\mid n$; if $E\ne D$, then
$d<n$ since $r=1$ would force $D=Ev_1$ and, as $1\in D$, $v_1\in E$.

\KEntry{K07.23.finite_division_ring}{Wedderburn's finite-division-ring theorem}
\begin{theorem}
Every finite division ring is commutative, and hence is a finite field.
\end{theorem}
\begin{proof}
Let the center have $q$ elements and $|D|=q^n$. If $n=1$, $D=Z(D)$.
Suppose $n>1$. Apply the class equation to the multiplicative group
$D^\times$. Its central elements are exactly $Z(D)^\times$: commuting
with all nonzero elements is equivalent to commuting with all elements.
For a representative $x_i$ of a noncentral conjugacy class its group
centralizer is $C_D(x_i)^\times$. By
\ref{K07.23.centralizers}, $|C_D(x_i)|=q^{d_i}$ with
$d_i\mid n$ and $d_i<n$. Orbit--stabilizer therefore gives
\[
 q^n-1=(q-1)+\sum_i\frac{q^n-1}{q^{d_i}-1}.
\]
Use the integer cyclotomic factorization established in
\ref{K07.15.cyclotomic_polynomials}. Since $d_i$ is a proper divisor of
$n$, cancellation of the factors indexed by divisors of $d_i$ shows that
$\Phi_n(q)$ divides each displayed quotient, and it also divides
$q^n-1$. It follows that $\Phi_n(q)$ divides $q-1$.

On the other hand, in $\C$ the roots of $\Phi_n$ are the primitive
$n$th roots $\zeta$. None equals $1$ when $n>1$. Each has absolute value
one and
\[
 |q-\zeta|^2=(q-1)^2+2q(1-\operatorname{Re}\zeta)>(q-1)^2.
\]
The strict inequality holds because a unit complex number with real part
one is one. Hence
$|\Phi_n(q)|=\prod_\zeta |q-\zeta|>(q-1)^{\varphi(n)}\ge q-1$;
the last inequality uses $q\ge2$ and $\varphi(n)\ge1$.
A nonzero integer of that absolute value cannot divide the positive
integer $q-1$. This contradiction proves $n=1$.
\end{proof}
No classification of finite simple groups, primitive-prime-divisor theorem,
or prior commutativity of the centralizers is used.

\KEntry{K07.23.algebraically_closed_center}{Algebraic division rings over an algebraically closed central field}
\begin{lemma}
Let $K$ be algebraically closed and embedded in the center of a division
ring $D$. If each element of $D$ is algebraic over $K$, then $D=K$.
More generally, for any central field $K$ and algebraic $x\in D$, the
subring $K[x]$ is a finite field extension of $K$.
\end{lemma}
\begin{proof}
Evaluation $K[X]\to D$ is a ring homomorphism because the coefficients
are central. Its kernel is nonzero by algebraicity and prime because
$D$ has no zero divisors. Its monic generator $m_x$ is irreducible:
a nontrivial factorization would evaluate to a product zero and yield an
annihilator of smaller degree. Polynomial division proves
$K[x]\simeq K[X]/(m_x)$, a finite field extension of degree $\deg m_x$.
If $K$ is algebraically closed, $m_x$ is linear, so $x\in K$.
Applying this to every $x$ proves the first assertion.
\end{proof}
In particular this lemma concerns algebraicity over a \emph{central}
field. Merely containing a copy of an algebraically closed field does not
force a division ring to equal it.

\KEntry{K07.23.quaternion_algebra}{Hamilton's quaternion algebra and its norm}
Define $\mathbb H=\C\oplus\C j$ as a real vector space, with pair
multiplication
\[
 (z,w)(z',w')=(zz'-w\overline{w'},\;zw'+w\overline{z'}).
\]
The map
\[
 (z,w)\longmapsto
 \begin{pmatrix}z&w\\-\overline w&\overline z\end{pmatrix}
\]
is injective and preserves this multiplication, as entry-by-entry matrix
multiplication verifies. Thus the product is associative and distributive,
and $(1,0)$ is its identity. In the notation
$i=(\sqrt{-1},0)$, $j=(0,1)$, $k=(0,\sqrt{-1})$,
$i^2=j^2=k^2=-1$, $ij=k=-ji$. The four elements $1,i,j,k$ form a
real basis because the two complex coordinates are independent real
two-dimensional coordinates.

Put $\overline{(z,w)}=(\overline z,-w)$ and
$N(z,w)=|z|^2+|w|^2$. Direct multiplication on either side gives
$u\overline u=\overline u u=N(u)1$. Thus a nonzero $u$ has inverse
$\overline u/N(u)$, and $\mathbb H$ is a division ring. Its matrix
determinant is $N(u)$; multiplicativity of the determinant proves
$N(uv)=N(u)N(v)$. For $u=a+bi+cj+dk$ the norm is
$a^2+b^2+c^2+d^2$. Finally $ui=iu$ forces $c=d=0$, and then
$uj=ju$ forces $b=0$. Its center is exactly $\R$. This constructs
the algebra itself, rather than identifying it solely by a multiplication
table with unverified associativity.

\KEntry{K07.23.frobenius}{Frobenius's real algebraic division-ring theorem}
\begin{theorem}
Suppose $D$ is an associative division ring with a specified central
embedding $\R\hookrightarrow D$, and every element of $D$ is algebraic
over $\R$. Then $D$ is isomorphic as a real algebra to exactly one of
$\R$, $\C$, $\mathbb H$. In particular the result holds for every
finite-dimensional associative real division algebra.
\end{theorem}
\begin{proof}
The algebraic fundamental theorem of algebra proved in
Appendix~\ref{app:algebraic-fta} implies that an irreducible real
polynomial has degree one or two: factor it into complex linear factors;
each nonreal root and its conjugate give a real quadratic factor, whereas
a real root gives a real linear factor. Each element of $D$ therefore
has minimal polynomial of degree at most two by
\ref{K07.23.algebraically_closed_center}. If $D\ne\R$, choose
$x\notin\R$. Its monic quadratic is $X^2+bX+c$, with negative
discriminant since it has no real root. Setting
$i=(x+b/2)/\sqrt{c-b^2/4}$ gives $i^2=-1$ and a subfield
$C=\R(i)\simeq\C$.

Let $D_+$ be the centralizer of $i$. It is a division subring and $C$
is central in it. Every one of its elements is algebraic over $C$, since
an annihilating real polynomial is also a complex-coefficient polynomial.
The preceding lemma applied to the algebraically closed field $C$ proves
$D_+=C$.

The real-linear map $T(y)=iyi^{-1}$ satisfies $T^2=1$ because
$i^2=-1$ is central. Hence
\[
 D=D_+\oplus D_-,\qquad
 D_-=\{y:iy=-yi\},\qquad
 y=\tfrac12(y+T(y))+\tfrac12(y-T(y)).
\]
The intersection is zero since the characteristic is zero. If $D_-=0$,
then $D=C$. Otherwise choose nonzero $j\in D_-$. For $z=a+bi\in C$
we have $jz=\overline z j$. For any $y\in D_-$, the element
$yj^{-1}$ commutes with $i$: both $y$ and $j^{-1}$ anticommute with
$i$. It therefore belongs to $C$. Thus $D_-=Cj$.

The element $j^2$ commutes with $i$, so it belongs to $C$. It also
commutes with $j$, while $jz=\overline z j$ shows that a complex
scalar commuting with $j$ is real. Thus $j^2=r\in\R$ and $r\ne0$.
It cannot be positive: if $r=s^2$ with $s>0$, the central factorization
$(j-s)(j+s)=0$ in a division ring would give $j=\pm s$, contrary
to $j\in D_-\setminus\{0\}$. Consequently $r<0$. Replacing $j$
by $j/\sqrt{-r}$ gives $j^2=-1$ without changing $Cj$.

Every element has a unique form $z+w j$, and its multiplication is
exactly the pair multiplication in \ref{K07.23.quaternion_algebra}:
use $jz'=\overline{z'}j$ and $j^2=-1$ to expand the product.
The resulting bijective real-linear multiplicative map
$\mathbb H\to D$ proves the assertion. These three real algebras have
different dimensions $1,2,4$ and hence are pairwise nonisomorphic.
If $D$ was assumed finite-dimensional, every element is algebraic by
linear dependence of $1,x,x^2,\ldots$; thus it satisfies the theorem's
stated hypothesis.
\end{proof}
Associativity and the central embedding are essential. The theorem does
not classify nonassociative real division algebras such as the octonions.
The argument also explains why algebraicity alone already implies finite
dimension in this particular real setting.

\KEntry{K07.23.real_schur}{The real form of Schur's endomorphism alternative}
\begin{corollary}
For a nonzero finite-dimensional irreducible real representation $V$ of
any group $G$, the algebra $\End_G(V)$ is isomorphic as a real algebra
to $\R$, $\C$ or $\mathbb H$.
\end{corollary}
\begin{proof}
The kernel and image of a nonzero equivariant endomorphism are invariant.
Irreducibility gives zero kernel and full image, and its inverse is
equivariant by multiplying the intertwining equation by the inverse.
Thus the endomorphism algebra is a division ring; the identity is nonzero
because $V\ne0$. It is a real-linear subspace of the finite-dimensional
algebra $\End_\R(V)$, so it is finite-dimensional. Scalar real maps
are central, with an injective scalar embedding since $V\ne0$.
Apply \ref{K07.23.frobenius}.
\end{proof}
No finiteness of $G$ is needed in this corollary. It replaces the false
real-field conclusion that every irreducible commutant consists only of
real scalar maps; the scalar version of Schur's lemma required an
algebraically closed field in the earlier representation chapter.
